\documentclass{report}
\usepackage{amsmath,amssymb}
\setlength{\parindent}{0mm}
\setlength{\parskip}{1em}
\begin{document}
\begin{center}
\rule{6in}{1pt} \
{\large Rosemary Renaut \\
{\bf Regularization by Unbiased Predictive Risk Estimation with Truncated Spectrum for LSQR solution of Large Scale Gravity Inversion}}

School of Mathematical and Statistical Sciences \\ Arizona State University \\ Tempe \\ AZ 85287-1804
\\
{\tt renaut@asu.edu}\\
Saeed  Vatankhah\end{center}

Tikhonov regularization for projected solutions of large-scale ill-posed
problems is considered. The LSQR algorithm is used to project the problem
onto a subspace and regularization then applied to find a subspace
approximation to the full problem. Determination of the regularization
parameter using the method of unbiased predictive risk estimation (UPRE)
is considered and contrasted with the generalized cross validation and
discrepancy principle techniques. Examining the unbiased predictive risk
estimator for the projected problem, it is shown that the obtained
regularized parameter provides a good estimate for that to be used for
the full problem with the solution found on the projected space. The
connection between regularization for full and projected systems for the
discrepancy and generalized cross validation estimators is also discussed
and an argument for the weight parameter in the weighted generalized
cross validation approach is provided. All results are independent of
whether systems are over or underdetermined, the latter of which has not
been considered in discussions of regularization parameter estimation for
projected systems.

Numerical simulations for standard one dimensional test problems and two
dimensional data for both image restoration and tomographic image
reconstruction support the analysis and validate the techniques. The size
of the projected problem is found using an extension of a noise revealing
function for the projected problem. Furthermore, an iteratively
reweighted regularization (IRLS) approach for edge preserving
regularization is extended for projected systems, providing stabilization
of the solutions of the projected systems with respect to the
determination of the size of the projected subspace.

This work is extended for sparse inversion of the large scale gravity
problem with the $L_{1}$-type stabilizer using IRLS. The regularization
parameter for the IRLS problem is estimated using the UPRE extended for
the projected problem. Further analysis leads to an improvement of the
projected UPRE via analysis based on truncation of the projected
spectrum. Numerical results for synthetic examples and real field data
demonstrate the efficiency of the method.


\end{document}
